% ─────────────────────────────────────────────────────────────────────────────
% Appendix G: Hyperparameter Tables
% ─────────────────────────────────────────────────────────────────────────────

\section*{Appendix G: Hyperparameter Tables}
\label{app:hyperparameters}

All values are sourced directly from the W\&B run configs (CSV logs) unless
marked $^?$ (inferred from thesis text or figure annotations).
Entries marked $^\dagger$ are flagged in Section~\ref{sec:notes_confounds}.
\texttt{max\_ep} is the maximum epochs per task; \texttt{ep/task} is the
configured training epochs; \texttt{max\_dirs} is the gradient memory column
budget.

% ─────────────────────────────────────────────────────────────────────────────
\subsection*{G.1 \quad Standalone Target-Network Experiments}
\label{app:hyperparameters_standalone}

All standalone experiments share the following fixed settings:
\texttt{model = TargetNetwork}, \texttt{normalize = False},
\texttt{regulizer = False}, \texttt{damping = 0.01}, \texttt{batch\_size = 10},
$n = 5$ seeds.
Learning rates are taken from \citet{garg_fisher-orthogonal_2026} Table~1
and vary per method; they are reproduced in Table~\ref{tab:hp_lr} for
reference.
The first-task optimizer is SGD for all MNIST benchmarks and Adam for
CIFAR benchmarks, at $\eta_{\mathrm{first}} = 10^{-3}$ throughout.
For MNIST benchmarks this deviates from the published value of
$\eta_{\mathrm{first}} = 10^{-5}$, which fails to converge
(see Appendix~D, Figure~\ref{fig:sgd_threshold}).

\begin{table}[H]
\centering
\caption{%
    \textbf{Benchmark-level training constants.}
    Fisher samples marked \textit{full} use the entire training set
    (60\,000 for Permuted-MNIST; 12\,000 for Split-MNIST).
}
\label{tab:hp_benchmark}
\small
\begin{tabular}{lccccc}
\toprule
Benchmark & Tasks & ep/task & Fisher samp. & EWC $\lambda$ & Proj.\ $\lambda$ \\
\midrule
Permuted-MNIST    & 5  & 5  & \textit{full} (60\,000) & 10     & $10^{-2}$  \\
Split-MNIST MH    & 5  & 5  & \textit{full} (12\,000) & 400    & $5\!\times\!10^{-4}$ \\
Split-MNIST SH    & 5  & 5  & \textit{full} (12\,000) & 400    & $5\!\times\!10^{-4}$ \\
Split-CIFAR10 MH  & 5  & 5  & 1\,024                  & 50     & $10^{-3}$  \\
Split-CIFAR100 MH & 10 & 10 & 1\,024                  & 10$^\ddagger$ & $10^{-3}$ \\
\bottomrule
\end{tabular}
\end{table}
\noindent$^\ddagger$CIFAR-100 MH: $\lambda = 50$ produced collapsed runs and
is excluded; $\lambda = 10$ is used (see Figure~\ref{fig:standalone_cifar100}).

\begin{table}[H]
\centering
\caption{%
    \textbf{Per-method learning rates} taken from
    \citet{garg_fisher-orthogonal_2026} Table~1.
    \texttt{iFOPNG} uses the same values as \texttt{FOPNG}.
}
\label{tab:hp_lr}
\small
\resizebox{\linewidth}{!}{%
\begin{tabular}{lcccccccc}
\toprule
Benchmark & Adam & SGD & EWC & FNG & OGD & ONG & FOPNG & iFOPNG \\
\midrule
Permuted-MNIST    & $10^{-4}$ & $5\!\times\!10^{-3}$ & $10^{-4}$          & $10^{-3}$          & $5\!\times\!10^{-3}$ & $5\!\times\!10^{-3}$ & $10^{-4}$          & $10^{-4}$ \\
Split-MNIST MH/SH & $10^{-5}$ & $5\!\times\!10^{-4}$ & $5\!\times\!10^{-4}$ & $10^{-3}$          & $5\!\times\!10^{-4}$ & $5\!\times\!10^{-4}$ & $10^{-5}$          & $10^{-5}$ \\
Split-CIFAR10 MH  & $10^{-3}$ & $5\!\times\!10^{-2}$ & $10^{-3}$          & $10^{-2}$          & $5\!\times\!10^{-2}$ & $10^{-2}$          & $10^{-3}$          & $10^{-3}$ \\
Split-CIFAR100 MH & $10^{-4}$ & $5\!\times\!10^{-3}$ & $10^{-2}$          & $5\!\times\!10^{-3}$ & $10^{-2}$          & $10^{-2}$          & $5\!\times\!10^{-3}$ & $5\!\times\!10^{-3}$ \\
\bottomrule
\end{tabular}}
\end{table}

% ─────────────────────────────────────────────────────────────────────────────
\subsection*{G.2 \quad Hypernetwork Experiments}
\label{app:hyperparameters_hn}

All HN experiments: \texttt{model = HyperNetwork}, \texttt{regulizer = True},
$\beta = 0.1$ (\texttt{lam = 0.001}), $\eta = 10^{-3}$, $n = 5$ seeds.

\begin{table}[H]
\centering
\caption{%
    \textbf{Hypernetwork architecture and training parameters.}
    \texttt{c\_emb} = chunk embedding dim; \texttt{t\_emb} = task embedding dim;
    \texttt{max\_dir} = maximum gradient memory directions.
    $^\dagger$Exp~411: results extracted from the 35 completed runs at
    chunk\,=\,6\,000, $d_h$\,=\,32, \texttt{normalize\,=\,True};
    remaining runs are exploratory pilots (see G.6).
}
\label{tab:hp_hn}
\resizebox{\linewidth}{!}{%
\begin{tabular}{lcccccccccc}
\toprule
Role / Benchmark & Tasks & ep/task & max\_ep
  & First opt & $d_h$ & chunk & c\_emb & t\_emb & \makecell{grads\\/ task} & max\_dir \\
\midrule
CIFAR-10 HN (main, Exp~408)              & 5  & 10 & 50 & AdamW & 32 & 6\,000 & 16 & 16 & 80  & 1\,000 \\
CIFAR-100 HN (Exp~411)$^\dagger$         & 10 & 10 & 50 & AdamW & 32 & 6\,000 & 16 & 16 & 200 & 800    \\
\midrule
MNIST HN ($d_h{=}8$, suffocated, Exp~407) & 5 & 15 & 15 & AdamW & 8  & 64 & 32 & 32 & 100 & 5\,000 \\
\midrule
$d_h$ sweep --- $d_h = 4$               & 5  & 15 & 15 & AdamW & 4  & 64 & 32 & 32 & 80  & 5\,000 \\
$d_h$ sweep --- $d_h = 8$               & 5  & 15 & 15 & AdamW & 8  & 64 & 32 & 32 & 100 & 5\,000 \\
$d_h$ sweep --- $d_h = 16$              & 5  & 15 & 15 & AdamW & 16 & 64 & 32 & 32 & 80  & 5\,000 \\
\bottomrule
\end{tabular}}
\end{table}

% ─────────────────────────────────────────────────────────────────────────────
\subsection*{G.3 \quad Normalization Ablation}
\label{app:hyperparameters_norm}

% All three conditions share the same architecture:
% Split-CIFAR10, 5 tasks, \texttt{max\_ep = 50},
% chunk\,=\,6\,000, $d_h$\,=\,32, batch\,=\,32,
% AdamW first-task at $\eta = 10^{-3}$, \texttt{iFOPNG} only.

\begin{table}[H]
\centering
\caption{%
    \textbf{Normalization ablation conditions (Sub-RQ2).}
    All conditions: Split-CIFAR10 HN, chunk\,=\,6\,000, $d_h$\,=\,32,
    \texttt{iFOPNG} only, AdamW first-task at $\eta = 10^{-3}$.
    Exp~410 (gradient-only) contains pilot runs at chunk\,=\,256,
    $d_h$\,=\,64; only chunk\,=\,6\,000, $d_h$\,=\,32 runs are used.
}
\label{tab:hp_norm}
\resizebox{\linewidth}{!}{%
\begin{tabular}{lccccccc}
\toprule
Condition & Exp & \texttt{normalize} & QR on $G$ & Max-scale $\hat{F}$
  & Seeds & Accuracy & BWT \\
\midrule
No normalization   & 409 & \texttt{False} & \texttimes  & \texttimes  & 5 & 72.2\% & $-0.042$ \\
Gradient-only      & 410 & \texttt{True}  & \checkmark  & \texttimes  & 3 & 68.9\% & $-0.060$ \\
Full normalization & 408 & \texttt{True}  & \checkmark  & \checkmark  & 6 & 86.0\% & $-0.006$ \\
\bottomrule
\end{tabular}}
\end{table}

% ─────────────────────────────────────────────────────────────────────────────
\subsection*{G.4 \quad Sub-RQ4 and Compression Ablation}
\label{app:hyperparameters_rq4}

Both Sub-RQ4 conditions and all three compression ablation conditions use
20-task Permuted-MNIST, \texttt{TargetNetwork}, \texttt{normalize = False},
batch\,=\,10.

\begin{table}[H]
\centering
\caption{%
    \textbf{Sub-RQ4 accumulation parameters} ($n = 5$ seeds,
    seeds $\{42, 111, 811, 1234, 2137\}$, $\alpha = 0.5$ for EMA).
}
\label{tab:hp_rq4}
\resizebox{0.90\linewidth}{!}{%
\begin{tabular}{lccccccc}
\toprule
Operator & Tasks & ep/task & max\_ep & LR & $\lambda$ & \makecell{grads\\/ task} & \makecell{Fisher\\samp.} \\
\midrule
MAX (\texttt{iFOPNG})                      & 20 & 5 & 5 & $10^{-4}$ & $10^{-2}$ & 80 & 60\,000 \\
EMA (\texttt{iFOPNG\_ema}, $\alpha = 0.5$) & 20 & 5 & 5 & $10^{-4}$ & $10^{-2}$ & 80 & 60\,000 \\
\bottomrule
\end{tabular}}
\end{table}

\begin{table}[H]
\centering
\caption{%
    \textbf{Compression ablation parameters} ($n = 5$ seeds each,
    seeds $\{42, 111, 811, 1234, 2137\}$).
    $\lambda = 10^{-3}$ differs from the Sub-RQ4 baseline ($\lambda = 10^{-2}$);
    the ablation is internally controlled but not directly comparable to
    Sub-RQ4 on this hyperparameter.
    $\alpha = 0.3$ for the internal F\_old EMA update in all three conditions,
    differing from the Sub-RQ4 value of $\alpha = 0.5$; results are therefore
    not directly comparable to Sub-RQ4 on this parameter either.
}
\label{tab:hp_compression}
\resizebox{0.90\linewidth}{!}{%
\begin{tabular}{lccccccc}
\toprule
Strategy & Tasks & ep/task & max\_ep & LR & $\lambda$ & \makecell{grads\\/ task} & max\_dirs \\
\midrule
SVD  (Exp~421) & 20 & 5 & 5 & $10^{-4}$ & $10^{-3}$ & 80 & 400 \\
FIFO (Exp~422) & 20 & 5 & 5 & $10^{-4}$ & $10^{-3}$ & 80 & 400 \\
STOP (Exp~423) & 20 & 5 & 5 & $10^{-4}$ & $10^{-3}$ & 80 & 400 \\
\bottomrule
\end{tabular}}
\end{table}

% ─────────────────────────────────────────────────────────────────────────────
\subsection*{G.5 \quad SGD Convergence Threshold Diagnostic}
\label{app:sgd_threshold}

The SGD threshold diagnostic (Figure~\ref{fig:sgd_threshold}) sweeps
$\eta \in \{10^{-5},\, 5\!\times\!10^{-5},\, 10^{-4},\, 2\!\times\!10^{-4},\,
5\!\times\!10^{-4},\, 10^{-3},\, 5\!\times\!10^{-3},\, 10^{-2}\}$
on Split-MNIST SH (\texttt{TargetNetwork}, 1 task, 5 epochs).
Convergence threshold: $\eta = 10^{-4}$.
All standalone experiments in this thesis use $\eta_{\mathrm{first}} = 10^{-3}$,
deviating from the \citet{garg_fisher-orthogonal_2026} protocol.